\documentclass[10pt,a4paper]{amsart}
\usepackage[margin=19mm]{geometry}
\usepackage{amsmath,amssymb,mathtools}
\usepackage[scr=rsfs,cal=euler]{mathalfa}
\usepackage{xfrac}
\usepackage[unicode,hidelinks,bookmarks=false]{hyperref}
\newcommand{\ie}{i.e.\ }
\newcommand{\cf}{cf.\ }
\newcommand{\resp}{respectively\ }
\newcommand{\Subsection}{\subsection}
\providecommand{\nobreakdash}{\nobreak\hbox{-}}
\setlength{\emergencystretch}{2em}
\multlinegap0pt
\usepackage{amssymb}
\usepackage{tikz}
\newtheorem{proposition}{Proposition}
\title{Extended VC-dimension, and Radon and Tverberg type theorems for unions of convex sets}
\author{Noga Alon, Shakhar Smorodinsky}
\date{Source: arXiv:2506.17777v3 (CC BY 4.0)}
\begin{document}
\maketitle

\noindent\emph{Source and licence.} Extended VC-dimension, and Radon and Tverberg type theorems for unions of convex sets. Version arXiv:2506.17777v3 (CC BY 4.0), distributed by arXiv under the Creative Commons Attribution 4.0 International licence, http://creativecommons.org/licenses/by/4.0/, as recorded in the arXiv licence metadata for this exact version and shown on the abstract page https://arxiv.org/abs/2506.17777v3. Full licence notice: \texttt{licenses/arxiv-2506.17777-CC-BY-4.0.txt}.\par\smallskip

\noindent\textit{Editorial scope.} Counted unit: the article's proposition t42 (source0.tex lines 876--883). The article's complete original proof of it is delivered with it (source0.tex lines 885--957, one \texttt{proof} environment of 73 lines). No mathematics is written, repaired or re-derived: the statement and the proof are the article's own lines, in the article's own notation, and no retained line is substituted or re-flowed.  No further source-local context is needed: every cross-reference of every retained line resolves inside the two delivered spans.  Nothing was omitted from any retained span, so the package carries no omission record.  No retained line contains a citation, so the wrapper carries no bibliography and no bibliography entry.  No retained line contains a figure environment, an image-inclusion command or drawing code, so no image is delivered and the package declares no asset dependency.  Editorial formatting only, in addition: an AMS \texttt{amsart} preamble that restates the article's own declarations and macros used by the retained lines, namely the environments \texttt{proposition} on separate counters, together with the standard packages those lines need.  The retained environments print in this document as Proposition 1 in order of appearance, whereas the article prints the same items with the numbering of its own full document.  The complete native source of the article as distributed in the arXiv e-print for this exact version is bundled unchanged as \texttt{sources/180272/source0.tex}.
\begin{proposition}
\label{t42}
For every $d \geq 1, s \geq 3$ and even $r \geq 2$,
$$
f_r(d,s) > \frac{1}{4} \left ( \left \lfloor \frac{d}{2}\right \rfloor+1\right)
\left \lfloor \frac{s-1}{2} \right \rfloor r^2
$$
\end{proposition}

\begin{proof}
Put $m=(\lfloor \frac{d}{2}\rfloor+1)r/2$,
$p=\lfloor \frac{s-1}{2} \rfloor r/2$ and $n=mp$.
Let $P$ be a set of $n$ points on the moment curve
in  $\mathbb{R}^d$ consisting of the $n$ points
$z_i=(t_i,t_i^2, \ldots t_i^d)$, where $0 <t_1<t_2 <  \ldots
<t_n<1$. Let  $I_1,I_2, \ldots ,I_p$ be $p$ open intervals that cover
$(0,1)$ and appear on it in this order, 
where the left endpoint of each interval $I_{j+1}$ is
just slightly smaller than the right endpoint of $I_j$ for each $j$.
The intervals are chosen so that each $t_i$ belongs to exactly one
such interval, and each interval contains exactly $m$ points $t_j$.

In order to complete the proof we show
that for any coloring of the points of $P$  by $r$ colors
there are $s$-convex sets $C_i$, with $C_i$ containing all points
of color $i$, so that the intersection of all the sets $C_i$ is
empty. Each set $C_i$  will be defined as the union of at most
$s$ convex sets, where each of the convex sets will be 
the convex hull of points of color $i$ with first coordinate 
lying in a union of some consecutive intervals $I_q$. The crucial
property of the definition of these convex sets is that for 
each interval $I_q$  there will be an index $i$ so that one of the
convex sets in the definition of $C_i$ will be the set of all
points $z_j=(t_j,t_j^2, \ldots ,t_j^d)$
of color $i$ for which $t_j$ lies in this single interval $I_q$.
Moreover, this will be done in a way that ensures that the number of such
points is always at most $\lfloor d/2 \rfloor$.  

Note that if such
a choice indeed exists, then the intersection of the corresponding 
$r$ sets $C_i$ will be empty as needed. Indeed, otherwise the
intersection is some point $z \in \mathbb{R}^d$ (not necessarily on
the moment curve). Let $t$ denote the first coordinate of this
point $z$. Suppose first that $t$ 
belongs only to one interval $I_q$, and let
$i$ be the color chosen for $I_q$ 
so that there are at most $\lfloor d/2 \rfloor$
points of color $i$ with their first coordinate in  $I_q$. 
Then $z$ has to lie in the convex hull of these points (as any
other convex set among the ones defining 
$C_i$ either has all its points with
first coordinate smaller than $t$ or all its points with 
first coordinate larger than $t$). Since any finite subset of the moment curve is
$\lfloor d/2 \rfloor$-neighborly\footnote{A set of points is $r$-neighborly if any subset of at most $r$ points form a face of the convex-hull of the set.}, this convex hull is disjoint from
the convex hull of all the other points of $P$, implying that
$z$ cannot lie in any other set $C_g$ besides $C_i$. If $t$ belongs
to two intervals, say $I_{q-1}$ and $I_q$  and $i$ is defined as
before, then $z$ cannot lie in $C_i$, since in this case any convex
set among those defining $C_i$
either has all its points with
first coordinate smaller than $t$ or all its points with 
first coordinate larger than $t$. 

It thus only remains to show that we can choose for every interval
$I_q$ a color $i$ with the required properties, making sure that no
fixed color is chosen more than $\lfloor (s-1)/2 \rfloor$ times
(as this way each set $C_i$ will be the union of at most $s$ convex 
sets). To do so we go over the intervals $I_q$ one by one, in an
arbitrary order (for example, from left to right). When dealing
with the interval $I_q$ after handling several previous ones, we
first note that since there are exactly 
$m=(\lfloor d/2 \rfloor+1)r/2$ 
points with first coordinate in $I_q$, there are at
least $r/2$ colors that appears at most 
$\lfloor d/2 \rfloor$ times in $I_q$. As so far we have selected the
colors $i$ for at most $p-1=\lfloor(s-1)/2 \rfloor r/2 -1$
intervals, there are less than $r/2$ colors $i$ that have already
been chosen $\lfloor(s-1)/2 \rfloor $ times. It follows that
there is at least one color $i$ that can be chosen for the interval
$I_q$, implying that the process terminates successfully. This
completes the proof of the theorem.
\end{proof}

\end{document}
